% =============================================================================
%  REPORT_PAPER_LNCS.tex — Conference paper (Springer LNCS/LNAI format)
%
%  Target venues: MIWAI 2026 (LNAI), SoICT 2026 (CCIS)
%
%  Compile:
%    cd /Users/kienha/spinet-v2/paper
%    pdflatex REPORT_PAPER_LNCS.tex
%    pdflatex REPORT_PAPER_LNCS.tex
%
%  Self-contained: no external .bib needed (uses thebibliography).
% =============================================================================

\documentclass[runningheads]{llncs}

\usepackage[T1]{fontenc}
\usepackage[utf8]{inputenc}
\usepackage{graphicx}
\usepackage{amsmath,amssymb}
\usepackage{booktabs}
\usepackage{multirow}
\usepackage{array}
\usepackage{microtype}
\usepackage{xcolor}
\usepackage{xspace}
\usepackage{enumitem}
\usepackage[colorlinks=true,linkcolor=blue,citecolor=blue,urlcolor=blue]{hyperref}

\graphicspath{{../experiments/v3_20260503/figures/}{../experiments/}{../experiments/paper_results/spider_zeroshot/}}

\definecolor{best}{HTML}{1A4DAB}
\newcommand{\best}[1]{\textcolor{best}{\textbf{#1}}}

% Compact float spacing
\setlength{\textfloatsep}{6pt}
\setlength{\intextsep}{6pt}
\setlength{\dbltextfloatsep}{6pt}
\setlength{\abovedisplayskip}{4pt}
\setlength{\belowdisplayskip}{4pt}

\begin{document}

\title{Hybrid 3D-CBAM Attention with Frozen Vision-Language
Embeddings for Lumbar Intervertebral Disc Grading
with Cross-Dataset Zero-Shot Transfer}

\titlerunning{Hybrid CBAM-BiomedCLIP for Lumbar IVD Grading}

\author{Trung Kien Ha\inst{1}\orcidID{0000-0000-0000-0000} \and
        Nh\^an Phan\inst{1}\orcidID{0000-0000-0000-0000}}

\authorrunning{T. K. Ha and N. Phan}

\institute{Faculty of Computer Science and Engineering,\\
           Ho Chi Minh City University of Technology (HCMUT), VNU-HCM,\\
           Ho Chi Minh City, Vietnam\\
           \email{\{kien.ha,nhanphan\}@hcmut.edu.vn}}

\maketitle

% =============================================================================
\begin{abstract}
Automated grading of lumbar intervertebral disc (IVD) pathologies on MRI is
a clinically valuable but technically difficult task: severity classes are
heavily imbalanced, scanner protocols vary across sites, and large-scale
labeled volumetric corpora are scarce. We present a two-branch hybrid
architecture that fuses a 3D ResNet-34 with Convolutional Block Attention
Modules (CBAM) and a frozen BiomedCLIP image-text encoder, trained on
the RSNA 2024 Lumbar Degenerative Classification dataset. The CBAM-3D
branch captures cross-slice anatomical context; the BiomedCLIP branch
contributes a semantic prior learned from 15 million biomedical
image-caption pairs. A trainable concat-MLP fuses the two streams into a
single 512-dimensional embedding aligned with text-prompt embeddings,
enabling both supervised classification and zero-shot label-space
extension. Combined with focal loss, square-root inverse class weights,
minority oversampling, and SpineNetV2-style augmentation, the hybrid
model improves Severe-class recall from 11.5\% to 46.4\% (4.0$\times$),
Severe F1 from 0.149 to 0.343 (+130\%), Severe AUC to 0.896, and Severe
AUPRC to 0.321 (7.6$\times$ over the class-prevalence random-baseline
AUPRC of 0.042) over a 3D ResNet-34 baseline. Cross-dataset
zero-shot transfer to SPIDER yields mean F1~$=$~0.362 across eight
unseen pathology labels using only text prompts. With lightweight
supervised transfer (15-epoch frozen-backbone heads, 3 random seeds),
Hybrid matches Baseline on aggregate SPIDER F1 ($0.653 \pm 0.014$
vs.\ $0.646 \pm 0.002$) while \emph{recovering above CBAM-only}
($0.619 \pm 0.010$, $\Delta = +0.034$, $> 2\sigma$) --- evidence that
the BiomedCLIP branch acts as a cross-dataset regularizer when CBAM
overfits RSNA-specific patterns. On clinical minority classes Hybrid
leads consistently: Disc Herniation recall from $0$ to $0.36 \pm
0.08$, Spondylolisthesis from $0.20$ to $0.30$, Pfirrmann Grade~4
from $0.47$ to $0.53$. To our knowledge no prior peer-reviewed
work has applied BiomedCLIP to RSNA 2024 lumbar classification at
the time of submission.
\end{abstract}

\keywords{Spine MRI \and Intervertebral disc \and Attention \and
Vision-language models \and BiomedCLIP \and CBAM \and Zero-shot
classification \and Class imbalance.}

% =============================================================================
\section{Introduction}
\label{sec:intro}

Low back pain is the leading cause of disability worldwide, with lumbar
disc degeneration driving most surgical and conservative interventions.
Quantitative MRI grading of intervertebral disc (IVD) pathology --- spinal
canal stenosis, left/right foraminal stenosis, Pfirrmann grade, and
related conditions --- remains predominantly manual and subject to
significant inter-rater variability~\cite{hallinan2021,nigru2024}.

Recent automated systems based on 3D convolutional networks
(SpineNet~\cite{jamaludin2017}, SpineNetV2~\cite{windsor2022}) have made
substantial progress, but two open problems persist:
\begin{enumerate}
\item \textbf{Class imbalance}: clinically the most important class
      (Severe) is typically less than 5\% of training data. Standard
      cross-entropy training yields models that ignore the minority
      class entirely (Severe Recall $\approx 0$\%).
\item \textbf{Limited label space}: a model trained on one schema (e.g.\
      RSNA 2024's 3-condition $\times$ 3-severity grid) cannot answer
      a query in another (e.g.\ SPIDER's Pfirrmann/spondylolisthesis/
      herniation schema) without retraining a fixed-class head.
\end{enumerate}

We address both problems with a single hybrid model. Our contributions,
ordered by strength of multi-seed evidence, are: (1) a
\textbf{cross-dataset regularization finding}: in a 3-seed SPIDER
transfer experiment, adding CBAM to the ResNet-34 backbone degrades
aggregate F1 by $-0.027$ (exceeds pooled seed $\sigma$), revealing
that the same attention module which helps in-domain RSNA Severe
metrics hurts cross-center generalization; the proposed Hybrid
(CBAM~+~frozen BiomedCLIP) recovers above CBAM-only by $+0.034$
($> 2\sigma$), establishing the BiomedCLIP branch as a regularizer.
(2) A \textbf{two-branch hybrid architecture} combining a CBAM-3D
ResNet-34~\cite{woo2018cbam,he2016resnet} with frozen
BiomedCLIP~\cite{zhang2023biomedclip} and a learned concat-MLP
fusion head. (3) \textbf{Cross-dataset label-space extension} via
cosine similarity against text-prompt embeddings, transferring to
SPIDER's eight unseen labels without fine-tuning (mean F1~=~0.362
zero-shot; Hybrid leads rare-class recall: Disc Herniation
$0.36 \pm 0.08$ vs.\ $0$ for Baseline/CBAM). (4) A
\textbf{class-imbalance handling pipeline} (focal
loss~\cite{lin2017focal} + square-root class weights + minority
oversampling + medical augmentation) lifting Severe-class recall
from 11.5\% to 46.4\% on RSNA 2024 (single-seed; 3-seed RSNA
evaluation pending, see \S\ref{sec:disc}).

To our knowledge no prior peer-reviewed work has applied BiomedCLIP to
RSNA 2024 lumbar classification, and no prior work has reported a
zero-shot label-space extension on the same architecture for lumbar
spine MRI.

% =============================================================================
\section{Related Work}
\label{sec:related}

\textbf{Lumbar disc grading.}
Jamaludin~et~al.~\cite{jamaludin2017} introduced SpineNet (VGG-style
2D); Windsor~et~al.~\cite{windsor2022} upgraded to a 3D ResNet-34
on per-IVD volumetric crops $(9,112,224)$, providing the
architectural foundation we adopt. Nigru~et~al.~\cite{nigru2024}
externally validated SpineNetV2 reporting foraminal balanced
accuracy $\approx 0.69$--$0.70$, which we treat as the sagittal-only
literature ceiling.

\textbf{Attention for medical imaging.}
The CBAM module~\cite{woo2018cbam} has been widely adapted to medical
2D problems. Lin~et~al.~\cite{lin2024cbam} combined CBAM with
multi-head self-attention on a \emph{balanced binary single-level}
RSNA 2024 subset, achieving F1~=~0.945 in that simplified setting.
We extend CBAM into a fully 3D form and combine it with a
foundation-model branch for full multi-level 3-class grading under
heavy imbalance.

\textbf{Vision-language models for medical imaging.}
BiomedCLIP~\cite{zhang2023biomedclip} pretrained a CLIP-style
image-text model on PMC-15M (15M PubMed caption pairs). Recent
benchmarking~\cite{datascarcity2024} shows BiomedCLIP linear
probing dominates supervised training in the few-shot regime,
motivating its use as a frozen extractor on small clinical
datasets. RadCLIP~\cite{lu2024radclip} extends CLIP with a
slice-pooling adapter and is the closest architectural alternative;
native 3D MRI VLMs (Decipher-MR~\cite{yang2026deciphermr}) were not
publicly available during this work.

\textbf{Positioning against prior RSNA 2024 methods.}
Table~\ref{tab:sota} summarizes published RSNA 2024 methods. Direct
numerical comparison is precluded by inconsistent setups:
binarization collapses the clinically critical Moderate--Severe
boundary, single-level (L4--L5) evaluation ignores spatial
heterogeneity, balanced-subset construction eliminates the very
class imbalance that motivates the problem, and Kaggle entries
report only weighted log-loss on $\geq 10$-backbone ensembles. Our
setup is intentionally the more difficult variant: full 3-class
severity grading across all five IVD levels on the natively
imbalanced distribution, evaluated as a single model.

\begin{table}[t]
\centering
\caption{Reported metrics from prior work on RSNA 2024 or closely
related lumbar spine MRI grading. Setups are inconsistent across
columns; see footnotes. ``—'' = not reported in the cited form.}
\label{tab:sota}
\renewcommand{\arraystretch}{1.05}
\setlength{\tabcolsep}{4pt}
\footnotesize
\begin{tabular}{@{}lcccl@{}}
\toprule
Method & Yr & Data & F1\,/\,AUC & Setup\,$^\dagger$ \\
\midrule
Lin~et~al.~\cite{lin2024cbam}            & '24 & RSNA\,sub & \best{.945}\,/\,.97 & Bin., bal., 1-lvl \\
Liu~et~al.~\cite{aktan2025}              & '25 & RSNA    & \best{.957}\,/\,—   & Bin., bal., joint seg \\
M-SCAN~\cite{hong2025mscan}              & '25 & RSNA    & —\,/\,.971 (canal) & Multi-view (sag.\,+\,ax.) \\
Phaphuang.~\cite{phaphuangwittayakul2026}& '26 & RSNA+ext. & \best{.956}\,/\,.99 & Bin., ext.\,bal. \\
nshen7~\cite{nshen2024}                  & '24 & RSNA & Sev.\,F1 \best{.501} & 3-cls, public split \\
\midrule
\textbf{Ours} (Hybrid)                   & '26 & RSNA    & \textbf{.528\,/\,.837} (Sev.\,F1 \textbf{.343}) & \textbf{3-cls, full multi-lvl, single model} \\
\bottomrule
\end{tabular}

\vspace{2pt}
\par
{\raggedright\scriptsize
$^\dagger$Setup abbreviations: Bin.~=~binarized to 2-class;
bal.~=~class-balanced subset; 1-lvl~=~single level (typically L4--L5);
ext.~=~external dataset.\par}
\end{table}


% =============================================================================
\section{Method}
\label{sec:method}

\subsection{Per-IVD Volume Extraction}

Given a sagittal T2 lumbar scan, we extract one volumetric region of
interest per intervertebral disc (L1--L2 through L5--S1) using the
keypoint annotations supplied with RSNA 2024. Each region is a tensor
$V \in \mathbb{R}^{1 \times 9 \times 112 \times 224}$ centered at the
disc midpoint, covering the full anatomical thickness of a single IVD.
Pixel intensities are clipped to the 1\textsuperscript{st}--99\textsuperscript{th}
percentile range and rescaled to $[0, 1]$. We follow the SpineNetV2
\cite{windsor2022} convention so that pretrained 3D ResNet-34 weights
can be loaded without modification.

\subsection{Hybrid Two-Branch Architecture}

Figure~\ref{fig:arch} illustrates the architecture. Two branches encode
the same per-IVD volume in complementary ways.

\textbf{CBAM-3D ResNet-34 branch (trainable).}
A 3D ResNet-34 with $3\!\times\!3\!\times\!3$ convolutions and BatchNorm3D
extracts volumetric features. After each of the four ResNet stages, given
an intermediate feature map
$F \in \mathbb{R}^{C \times D \times H \times W}$, a CBAM block
applies channel attention $M_c(F) = \sigma(\text{MLP}(\text{AvgPool}(F))
+ \text{MLP}(\text{MaxPool}(F)))$ (bottleneck ratio $r=16$),
followed by spatial attention $M_s(F') = \sigma(f^{7 \times 7
\times 7}([\text{AvgPool}_c(F'); \text{MaxPool}_c(F')]))$ using a
3D convolution. The composed refinement is $F'' = M_s(M_c(F) \otimes
F) \otimes (M_c(F) \otimes F)$, with $\otimes$ element-wise
multiplication. After global average pooling the branch outputs
$f_{\text{cbam}} \in \mathbb{R}^{512}$.

\textbf{BiomedCLIP branch (frozen, 2.5D).}
We split the volume into its 9 sagittal slices
$\{s_i\}_{i=1}^{9}$, replicate each slice to 3 channels, resize to
$224\!\times\!224$, and encode independently with frozen BiomedCLIP
ViT-B/16, $g_\theta : \mathbb{R}^{3\times 224 \times 224} \to
\mathbb{R}^{512}$. The 9 slice embeddings $z_i = g_\theta(s_i)$ are
aggregated by a learned attention pool:
\begin{equation}
\alpha_i = \frac{\exp(w_a^\top z_i)}{\sum_{j=1}^{9} \exp(w_a^\top z_j)},
\quad
f_{\text{bmc}} = \sum_{i=1}^{9} \alpha_i \, z_i,
\label{eq:slicepool}
\end{equation}
where $w_a \in \mathbb{R}^{512}$ is the only trainable parameter of
this branch. Replacing mean pooling with this attention pool allows the
branch to up-weight slices in which the lesion is most clearly visible.

\textbf{Fusion head (trainable).}
The two embeddings are concatenated and projected through a two-layer
MLP: $f_{\text{img}} = W_2\,\phi(W_1\,[f_{\text{cbam}};f_{\text{bmc}}]
+ b_1) + b_2$ with $W_1 \in \mathbb{R}^{768 \times 1024}$,
$W_2 \in \mathbb{R}^{512 \times 768}$, ReLU $\phi$, dropout $0.1$,
followed by L2 normalization $\hat{f}_{\text{img}} = f_{\text{img}}
/ \|f_{\text{img}}\|_2$.

\textbf{Classification heads.}
For RSNA we attach three task-specific 3-class linear heads (canal,
left foraminal, right foraminal). For zero-shot transfer we replace
each head with cosine similarity against text-prompt embeddings
$\{t_k\}$ produced by the BiomedCLIP text encoder, predicting
$\hat{y} = \arg\max_{k} s \cdot \hat{f}_{\text{img}}^{\top} t_k$
where $s$ is the BiomedCLIP-pretrained temperature. Changing the
prompt vocabulary changes the label space without retraining.

The trainable parameter budget is approximately 1.18\,M (CBAM-3D
fine-tuning + slice attention pool + fusion MLP), versus 218\,M
total parameters once the frozen backbones are counted.

\begin{figure}[t]
\centering
\includegraphics[width=0.98\linewidth]{hybrid_architecture.png}
\caption{Two-branch hybrid architecture. CBAM-3D learns RSNA-specific
volumetric features; frozen BiomedCLIP provides a semantic prior. The
fusion MLP and slice attention pool ($\sim$1.18\,M trainable parameters)
are the only learned components, and the cosine-similarity head allows
both supervised and zero-shot classification.}
\label{fig:arch}
\end{figure}

\subsection{Training Objective}

We address class imbalance with a four-component pipeline. The
per-task loss is the focal cross-entropy~\cite{lin2017focal}
$\mathrm{FL}_t(p_t) = -w_{c(t)}(1 - p_t)^{\gamma} \log(p_t)$ with
$\gamma = 2.0$, summed over the three RSNA conditions, with missing
labels excluded via \verb|ignore_index = -1|. Class imbalance is
further addressed by (i) \emph{square-root inverse class weights}
$w_c = (1/\sqrt{n_c})/Z$ with $Z$ normalizing the mean to 1, which
reduces the Severe/Normal gradient ratio from $\sim15\times$
(linear $1/n_c$) to $\sim 4\times$; (ii) \emph{minority oversampling}
factor 3--5 at the sampler stage, complementary to $w_c$ as it
changes batch composition rather than per-sample gradient; and
(iii) \emph{medical augmentation} following
SpineNetV2~\cite{windsor2022} \S4.4.3 (horizontal flip with
left/right label swap, rotation $\pm 10^{\circ}$, brightness/contrast
$\pm 20\%$, Gaussian noise $\sigma = 0.05$ at probability $0.3$).
Training uses AdamW, lr $10^{-4}$, batch size 32, weight decay
$10^{-4}$, up to 20 epochs with patience-5 early stopping on a
single NVIDIA RTX 4090 (16\,GB peak memory).

% =============================================================================
\section{Experiments}
\label{sec:exp}

\subsection{Datasets}

\textbf{RSNA 2024 LumbarDISC.} Approximately 2{,}000 patients with
sagittal T2 stacks. We use a patient-level 80/20 train/validation split
(\verb|random_state=42|), yielding 1{,}942 validation IVDs from 395
patients. Labels span 5 levels $\times$ 5 conditions $\times$ 3
severities. We restrict evaluation to three of the five conditions
(canal stenosis, left and right neural foraminal narrowing) for which
sagittal-only input is most clinically appropriate; left and right
subarticular stenosis, which are best assessed on axial reformats,
are left to future work and discussed in \S\ref{sec:disc}. Class
distribution per evaluated condition averages Normal/Mild $\approx 85\%$,
Moderate $\approx 10\%$, Severe $\approx 5\%$.

\textbf{SPIDER.}~\cite{vandergraaf2024spider} Approximately 1{,}400
IVDs from 218 patients across four Dutch hospitals (T1 and T2-FS
sagittal sequences). For consistency with the RSNA-trained encoder we
use the T2-FS sagittal series only. The label schema differs from
RSNA: Pfirrmann grades 1--5~\cite{pfirrmann2001}, Modic
changes 0--3~\cite{modic1988},
spondylolisthesis (binary), disc herniation (binary), disc narrowing
(binary), disc bulging (binary), upper-/lower-endplate damage
(binary). Used for cross-dataset zero-shot evaluation only; no SPIDER
images are seen during RSNA training.

\subsection{Evaluation Metrics}

We report macro-averaged Recall, Precision, F1, AUC (one-vs-rest), and
AUPRC, plus per-class Severe metrics. With $TP$, $FP$, $FN$ denoting
true positives, false positives, and false negatives at a threshold,
the threshold-dependent metrics are
\begin{equation}
\mathrm{Recall} = \frac{TP}{TP + FN},\quad
\mathrm{Precision} = \frac{TP}{TP + FP},
\label{eq:recprec}
\end{equation}
\begin{equation}
\mathrm{F1} = 2\cdot\frac{\mathrm{Precision}\cdot\mathrm{Recall}}
                         {\mathrm{Precision} + \mathrm{Recall}}.
\label{eq:f1}
\end{equation}
The threshold-free metrics integrate the ROC and PR curves:
\begin{equation}
\mathrm{AUC} = \int_{0}^{1} \mathrm{TPR}\!\big(\mathrm{FPR}^{-1}(u)\big)\,du,
\;\;
\mathrm{AUPRC} = \int_{0}^{1} P(R)\,dR.
\label{eq:auc_auprc}
\end{equation}
Macro averaging across the $C{=}3$ severities, $\frac{1}{C}\sum_{c}
\mathrm{metric}_c$, prevents the dominant Normal/Mild class (85\%
prevalence) from concealing failure on the minority Severe class. AUPRC
is reported alongside AUC because, on heavily imbalanced data, AUC can
remain high while a model fails to rank Severe cases meaningfully; the
AUPRC random-baseline equals the positive-class prevalence
($\approx 0.05$ for Severe).

% =============================================================================
\section{Results}
\label{sec:results}

\subsection{Main Results: Ablation Across Four Configurations}

Table~\ref{tab:main} reports nine metrics on the RSNA 2024 validation
split for four configurations: a 3D ResNet-34 baseline, CBAM-only
(no BiomedCLIP branch), BMC-only (no CBAM branch), and the full
hybrid. The hybrid wins 8/9 metrics; the only metric on which it does
not win is Mean Accuracy, which we discuss below. We note that the
nshen7~\cite{nshen2024} Severe F1 of $0.501$ in Table~\ref{tab:sota}
appears higher than our $0.343$, but reflects a different evaluation
protocol (alternative per-condition aggregation and a different
validation split) reported only as a public GitHub artifact without
peer review; our
$0.343$ is the macro-averaged Severe F1 over all three RSNA
conditions (canal, left foraminal, right foraminal), where the two
foraminal conditions are known to exhibit a sagittal-only ceiling
near $0.18$--$0.20$ (\S\ref{sec:disc}). A like-for-like comparison
under matched evaluation protocol is not currently possible from
public artifacts.

\begin{table}[t]
\centering
\caption{Ablation on RSNA 2024 validation (1{,}942 IVDs / 395 patients).
\textbf{Bold} = best per row.}
\label{tab:main}
\renewcommand{\arraystretch}{1.05}
\begin{tabular}{@{}lcccc@{}}
\toprule
Metric & Base & CBAM-only & BMC-only & \best{Hybrid} \\
\midrule
Mean Accuracy        & \best{81.4\%} & 68.98\% & 69.19\% & 72.1\% \\
Mean F1 macro        & 0.420 & 0.502 & 0.464 & \best{0.528} \\
Mean Recall macro    & 0.411 & 0.569 & 0.528 & \best{0.592} \\
Mean Precision macro & 0.486 & 0.510 & 0.475 & \best{0.517} \\
Mean AUC macro       & 0.826 & 0.822 & 0.812 & \best{0.837} \\
Mean AUPRC macro     & 0.523 & 0.519 & 0.482 & \best{0.526} \\
\midrule
\textbf{Severe F1}    & 0.149 & 0.304 & 0.229 & \best{0.343} \\
\textbf{Severe AUC}   & 0.860 & 0.886 & 0.874 & \best{0.896} \\
\textbf{Severe AUPRC} & 0.276 & 0.319 & 0.218 & \best{0.321} \\
\bottomrule
\end{tabular}
\end{table}

\textbf{Hybrid $>$ either branch alone.} CBAM-only achieves F1~=~0.502,
BMC-only achieves F1~=~0.464; the hybrid achieves 0.528. Each branch
addresses a blind spot of the other: CBAM-3D provides volumetric and
RSNA-specific features that the frozen 2D BiomedCLIP cannot learn,
while BiomedCLIP provides a semantic prior from 15M PubMed Central
pairs that the small RSNA training set cannot match. The fusion gain
(+2.6 percentage points over the better single branch) confirms the
two information sources are complementary rather than redundant.

\textbf{Severe class is the headline result.} Severe F1 increases
2.3$\times$ from 0.149 (baseline) to 0.343 (hybrid), Severe Recall
4.0$\times$ from 11.5\% to 46.4\%, and Severe AUC from 0.860 to 0.896.
Severe AUPRC of 0.321 against the class-prevalence random-baseline AUPRC
of 0.042 (Severe support averages 4.20\% across the three conditions)
represents a 7.6$\times$ gain over chance on the rare class.

\textbf{Why Mean Accuracy decreases.} The 9.3 percentage-point
accuracy drop is a direct consequence of rebalancing: the baseline's
high accuracy reflects a model that predicts Normal/Mild on most cases
(Severe Recall $= 11.5\%$). The hybrid trades majority-class accuracy
for minority-class recall, which is the clinically appropriate
trade-off in a screening setting where missing a Severe case is more
costly than a false alert.

\subsection{Cross-Dataset Zero-Shot Transfer to SPIDER}

We freeze the RSNA-trained Hybrid and evaluate on SPIDER directly,
replacing the classification heads with cosine similarity against
text prompts of the eight SPIDER label categories (no SPIDER
training data used). Table~\ref{tab:spider} reports per-disease
zero-shot F1; mean F1 across the eight unseen labels is $0.362$,
which is structurally impossible for Baseline and CBAM-only because
their fixed 3-class heads cannot accommodate the new label space
without retraining.

\begin{table}[t]
\centering
\caption{Zero-shot transfer of RSNA-trained Hybrid to SPIDER (no SPIDER
fine-tuning). Disc-related labels (top group) transfer strongly;
non-disc labels (bottom group) show a clear semantic gap.}
\label{tab:spider}
\renewcommand{\arraystretch}{1.05}
\begin{tabular}{@{}lcc@{}}
\toprule
SPIDER label              & F1     & AUC \\
\midrule
Disc bulging              & \best{0.709} & 0.806 \\
Disc herniation           & 0.522 & 0.802 \\
Disc narrowing            & 0.388 & 0.836 \\
\midrule
Lower endplate            & 0.452 & 0.747 \\
Upper endplate            & 0.441 & 0.743 \\
Pfirrmann (5-class)       & 0.147 & --     \\
Modic (4-class)           & 0.117 & --     \\
Spondylolisthesis         & 0.030 & 0.745 \\
\midrule
\textbf{Mean (all 8)}     & \best{0.362} & --     \\
\bottomrule
\end{tabular}
\end{table}

The selective-transfer pattern is itself a finding: disc-pathology
labels (bulging, herniation, narrowing) transfer well while
semantically distant labels (Modic, Pfirrmann, spondylolisthesis)
show a measurable domain gap, suggesting text-prompt zero-shot
inherits the semantic proximity of the contrastive pretraining
corpus. A Naked-BiomedCLIP baseline (no CBAM, no RSNA training)
underperforms our Hybrid on the same disc-related labels,
confirming that CBAM and the RSNA-supervised fusion head contribute
features that transfer beyond the original training domain.

\subsection{Cross-Dataset Transfer Learning on SPIDER 8-Label Schema}
\label{sec:spider_transfer}

To assess the full transfer learning potential beyond zero-shot, we
fine-tune lightweight task-specific heads (RSNA backbone frozen) on
the SPIDER training split for 15 epochs across all four configurations
and the complete 8-label schema (Pfirrmann 5-class, Modic 4-class, six
binary labels). The validation split has 234 IVDs from 34 patients;
backbone parameters are frozen for fair comparison across the four
configurations; only condition-specific heads (and the fusion MLP for
Hybrid/BMC-only) are trained. This is the natural complement to the
zero-shot evaluation above: zero-shot tests emergent text-prompt
extension, while supervised transfer tests how well each pretrained
representation maps to a different label space with minimal SPIDER
supervision.

\begin{table}[t]
\centering
\caption{SPIDER transfer learning, 4 configurations, mean across 8
conditions. All metrics on val (234 IVDs), reported as mean $\pm$
standard deviation across 3 random seeds ($\{42, 123, 456\}$).
\textbf{Bold} = best per row when the gap to second place exceeds
$2\sigma$ (the pooled standard deviation).}
\label{tab:spider_transfer_agg}
\renewcommand{\arraystretch}{1.05}
\footnotesize
\begin{tabular}{@{}lcccc@{}}
\toprule
Metric & Base & CBAM & BMC-only & Hybrid \\
\midrule
Mean Accuracy        & \best{80.23 $\pm$ 0.94} & 79.27 $\pm$ 1.19 & 76.12 $\pm$ 0.83 & 78.30 $\pm$ 1.64 \\
Mean F1 macro        & 0.646 $\pm$ 0.002 & 0.619 $\pm$ 0.010 & 0.613 $\pm$ 0.006 & 0.653 $\pm$ 0.014\,$^\ast$ \\
Mean AUC             & 0.842 $\pm$ 0.006 & 0.824 $\pm$ 0.018 & 0.808 $\pm$ 0.011 & 0.846 $\pm$ 0.018\,$^\ast$ \\
Mean AUPRC           & \best{0.633 $\pm$ 0.026} & 0.580 $\pm$ 0.016 & 0.544 $\pm$ 0.007 & 0.611 $\pm$ 0.010 \\
\midrule
Inference (ms/sa)    & 65.5 $\pm$ 2.7 & 66.7 $\pm$ 3.2 & 65.2 $\pm$ 3.9 & 66.4 $\pm$ 3.7 \\
\bottomrule
\end{tabular}

\vspace{2pt}
\par
{\raggedright\scriptsize
$^\ast$Hybrid and Baseline are within $2\sigma$ on Mean F1 / Mean AUC;
both pass-zone results. The Hybrid--CBAM gap on Mean F1 (Hybrid
$0.653$ vs.\ CBAM $0.619$, $\Delta = +0.034$, pooled $\sigma =
0.012$) and on Mean AUPRC ($\Delta = +0.031$, pooled $\sigma =
0.013$) exceeds $2\sigma$ and is therefore reported as a substantive
gap.\par}
\end{table}

\textbf{Aggregate result (3 seeds).} Three findings emerge from the
multi-seed analysis:

\begin{enumerate}[leftmargin=*,topsep=0pt]
\item \emph{Hybrid vs.\ Baseline}: the F1 gap of $+0.006$ is within
single-seed noise (pooled $\sigma = 0.010$). On Mean AUPRC, Baseline
actually leads slightly ($0.633$ vs.\ $0.611$). We do not claim
Hybrid surpasses Baseline on aggregate SPIDER metrics; instead, we
report that the two are at parity within seed variability, with the
relative ordering depending on the metric chosen.

\item \emph{Hybrid vs.\ CBAM}: this gap is significant and consistent
($+0.034$ F1, pooled $\sigma = 0.012$; $+0.031$ AUPRC, pooled
$\sigma = 0.013$). CBAM alone, optimized for RSNA's spatial patterns,
degrades on cross-dataset transfer; Hybrid recovers this loss by
fusing CBAM with BiomedCLIP's dataset-agnostic semantic prior.

\item \emph{CBAM hurts cross-dataset transfer}: CBAM Mean F1
($0.619$) is below Baseline ($0.646$) by $-0.027$, exceeding the
pooled CBAM--Baseline noise ($\sigma_{\text{pooled}} = 0.007$,
$|\Delta|/\sigma = 3.8$). This is a substantive
observation: the same CBAM module that improves in-domain (RSNA)
Severe-class metrics produces a measurable regression when the
backbone is reused on a different dataset. We discuss the implications
in \S\ref{sec:disc}.
\end{enumerate}

\textbf{Per-condition F1 and rare-class recovery.}
Table~\ref{tab:spider_minority} highlights conditions with rare
positive class ($\leq 5\%$ prevalence). For Spondylolisthesis
(positive prevalence $2.8\%$), Hybrid recovers $30.0\% \pm 9.8$ of
positive cases vs.\ Baseline $20.0\% \pm 9.8$, CBAM $0.0\%$, and
BMC-only $5.6\% \pm 7.9$. For Disc Herniation (positive prevalence
$4.7\%$), Hybrid achieves $35.6\% \pm 7.9$ recall while Baseline and
CBAM both fail to recall any positive cases ($0.0\%$). We highlight
that BMC-only (BiomedCLIP branch without CBAM) already achieves
$27.0\% \pm 15.4$ recall on this label; the BiomedCLIP semantic
prior is responsible for the bulk of the herniation recovery, while
the CBAM volumetric branch contributes an incremental
$+0.086~\pm~17.3$ when fused into the Hybrid. The two branches
therefore play complementary roles: BiomedCLIP enables minority-class
recovery via dataset-agnostic semantic features, and CBAM contributes
volumetric refinement on top. For Pfirrmann
Grade~4, Hybrid leads with $0.534 \pm 0.043$ vs.\ Baseline
$0.473 \pm 0.062$. Modic Types~1 and 3 (population prevalence $0.3\%$
and $0.5\%$ respectively) remain at $0\%$ recall across all four
configurations and are not amenable to 15-epoch frozen-backbone
training; they represent the lower end of what supervised transfer
can achieve on this dataset.

\begin{table}[t]
\centering
\caption{Minority-class recall on SPIDER (positive label for binary,
Grade~4 for Pfirrmann). Rare positive prevalence (\%) shown
parenthetically. Hybrid leads on rare/severe classes; aggregate F1
gains are concentrated in these cases.}
\label{tab:spider_minority}
\renewcommand{\arraystretch}{1.05}
\footnotesize
\begin{tabular}{@{}lcccc@{}}
\toprule
Condition (prevalence)              & Base   & CBAM   & BMC    & Hybrid \\
\midrule
Spondylolisthesis Yes (2.8\%)       & 0.200 $\pm$ 0.098 & 0.000 $\pm$ 0.000 & 0.056 $\pm$ 0.079 & \best{0.300 $\pm$ 0.098} \\
Disc Herniation Yes (4.7\%)         & 0.000 $\pm$ 0.000 & 0.000 $\pm$ 0.000 & 0.270 $\pm$ 0.154 & \best{0.356 $\pm$ 0.079} \\
Pfirrmann Grade~4 (19\%)            & 0.473 $\pm$ 0.062 & 0.396 $\pm$ 0.038 & 0.437 $\pm$ 0.163 & \best{0.534 $\pm$ 0.043} \\
Modic Type~1 (0.3\%)                & 0.000  & 0.000  & 0.000  & 0.000 \\
Modic Type~3 (0.5\%)                & 0.000  & 0.000  & 0.000  & 0.000 \\
\bottomrule
\end{tabular}
\end{table}

\textbf{Architectural advantage: zero-shot label extension.} Two of
the four configurations (BMC-only and Hybrid) accommodate new labels
at inference time by encoding a new text prompt through BiomedCLIP's
text encoder, requiring no retraining. Baseline and CBAM, which use
fixed linear heads, cannot extend their label space without retraining.
This is an architectural property independent of the F1 metric and
relevant for clinical deployment in centers that may need to
incorporate additional pathology labels (e.g.\ Schmorl's nodes,
Tarlov cysts) without additional supervised data collection.

\subsection{Per-condition Severe recall and qualitative analysis}

Per-condition Severe recall breakdown shows $70.4\%$ on canal
stenosis but only $18$--$22\%$ on left/right foraminal stenosis;
the foraminal recall floor is consistent with the sagittal-only
literature ceiling reported by Nigru~et~al.~\cite{nigru2024}
(foraminal balanced accuracy $\approx 0.69$--$0.70$). Grad-CAM
activation on Severe-class canal stenosis cases (computed against
the final CBAM stage and trilinearly upsampled) concentrates on the
central canal narrowing rather than diffusing across the disc,
providing qualitative support for the CBAM spatial-attention
hypothesis. Hybrid precision-recall curves for the Severe class
dominate the moderate-recall regime ($\mathrm{Recall} \in [0.2,
0.6]$), the operating range relevant to clinical triage.

\subsection{Computational Cost}

On a single RTX 4090, Hybrid training converges at epoch 10 in
24.8\,min (vs.\ Baseline 16.4\,min ep~18, CBAM 56.0\,min ep~14).
Inference latency is 9.65\,ms/sample (Hybrid) vs.\ 5.61\,ms
(Baseline), 5.60\,ms (CBAM), and 3.98\,ms (BMC-only); the Hybrid
overhead of $+72\%$ over CBAM-only stems from the frozen BiomedCLIP
forward pass. Throughput of 103.6 samples/s on the Hybrid is well
above clinical real-time requirements.

% =============================================================================
\section{Discussion}
\label{sec:disc}

\textbf{Why the hybrid works: complementary roles + cross-dataset
regularization.}
The two branches are complementary in information content:
CBAM-3D answers \emph{where} a lesion is by learning task-specific
volumetric features under direct supervision, while frozen
BiomedCLIP answers \emph{what} a lesion means by carrying a semantic
prior from 15M biomedical caption pairs. We use ``complementary''
in the empirical sense that each branch alone falls short of the
Hybrid; we do not claim a formal synergy effect, which would
require an interaction-term test (deferred). On in-domain RSNA,
CBAM is the single largest source of Severe-class improvement
($+0.155$ over Baseline). On out-of-domain SPIDER, however, CBAM
Mean F1 ($0.619 \pm 0.010$) \emph{drops below} the plain ResNet-34
Baseline ($0.646 \pm 0.002$) --- a $\Delta = -0.027$ regression
exceeding the pooled noise ($\sigma_{\text{pooled}} = 0.007$,
$|\Delta|/\sigma = 3.8$). The Hybrid ($0.653 \pm 0.014$) recovers
above CBAM by $+0.034$ ($> 2\sigma$) and matches Baseline within
seed variability. We frame Hybrid not as additive feature gain
beating Baseline, but as a regularizer that \emph{restores}
Baseline-level performance when CBAM is added: BiomedCLIP provides
a robust fallback when CBAM's in-domain bias is misaligned with
the target distribution. The same pattern explains why CBAM-only
fails on SPIDER Spondylolisthesis (Recall $0.000$) while Hybrid
recovers $0.300 \pm 0.098$, and why Disc Herniation recall goes
from $0$ for Baseline/CBAM to $0.356 \pm 0.079$ for Hybrid. A
trivial-ensemble alternative --- logit averaging of Baseline and
BMC-only --- cannot be fully ruled out from current experiments
and is deferred to future work.

\textbf{Clinical positioning and foraminal weakness.}
Severe Recall $46.4\%$ on RSNA, while a $4.0\times$ improvement
over Baseline, is below typical clinical screening sensitivity
thresholds (often $\geq 80\%$). At the current operating point the
model is best positioned as a \emph{second-reader assistance tool}
that flags candidate Severe cases for radiologist confirmation,
not a standalone screening system. The per-condition breakdown
reveals a structural limitation: canal stenosis Severe recall is
$70.4\%$ but left/right foraminal stenosis remains $18$--$22\%$.
Foraminal stenosis is most accurately assessed on \emph{axial}
reformats; we elected sagittal-only input to retain SpineNetV2
ResNet-34 weight compatibility and a consistent crop $(9, 112,
224)$. Nigru~et~al.~\cite{nigru2024} report foraminal balanced
accuracy $\approx 0.69$--$0.70$ on sagittal-only input as the
corresponding ceiling, and multi-view fusion (e.g.\ M-SCAN-style
sagittal--axial cross-attention) is the natural next step.

\textbf{Limitations and future work.}
(i) RSNA results are single-seed (42); SPIDER uses 3 seeds. A
3-seed RSNA sweep and bootstrap 95\% CI ($n=1000$) on the RSNA
Severe metrics are committed for the camera-ready version.
(ii) BiomedCLIP is frozen and 2D, processing slices independently;
native-3D MRI VLMs (Decipher-MR~\cite{yang2026deciphermr},
Merlin~\cite{blankemeier2026merlin}) were unavailable during this
work. (iii) The fusion strategy is restricted to concat-MLP; gated
fusion and cross-attention remain unexplored.
(iv) RSNA 2024 labels were produced via multi-annotator
crowdsourcing~\cite{kaggle2024} without released per-class
inter-rater statistics, so the achievable upper bound reflects
label noise as well as model capacity. Future work: RadCLIP
slice-pooling~\cite{lu2024radclip}, prompt-tuning via
BiomedCoOp~\cite{koleilat2025biomedcoop}, multi-view axial-T2
fusion, and external evaluation on the Mendeley Pfirrmann subset.

% =============================================================================
\section{Conclusion}
\label{sec:concl}
We presented a hybrid CBAM-3D + frozen BiomedCLIP architecture for
lumbar IVD grading on RSNA 2024. The model improves Severe-class
metrics (Recall $\times$4.0, F1 $\times$2.3, AUC 0.896) at a
controlled accuracy trade-off, and uniquely enables zero-shot
label-space extension to SPIDER's eight unseen labels via
text-prompt cosine similarity.

% =============================================================================
\begin{thebibliography}{99}
\footnotesize
\setlength{\itemsep}{1pt}

\bibitem{hallinan2021}
J.~T. P.~D. Hallinan \textit{et al.},
``Deep learning model for automated detection and classification of
central canal, lateral recess, and neural foraminal stenosis at lumbar
spine MRI,''
\textit{Radiology}, vol.~300, no.~1, pp.~130--138, 2021.

\bibitem{nigru2024}
S.~M. Nigru \textit{et al.},
``Automated grading of lumbar disc degeneration: external validation
on a multi-cohort dataset,''
\textit{North American Spine Society Journal}, 2024.

\bibitem{jamaludin2017}
A.~Jamaludin \textit{et al.},
``SpineNet: automated classification and evidence visualization in
spinal MRIs,''
\textit{Medical Image Analysis}, vol.~41, pp.~63--73, 2017.

\bibitem{windsor2022}
R.~Windsor \textit{et al.},
``SpineNetV2: automated detection, labelling and radiological grading
of clinical MR scans,''
\textit{arXiv:2205.01683}, 2022.

\bibitem{woo2018cbam}
S.~Woo, J.~Park, J.-Y. Lee, and I.~S. Kweon,
``CBAM: convolutional block attention module,''
in \textit{Proc. ECCV}, 2018, pp.~3--19.

\bibitem{he2016resnet}
K.~He, X.~Zhang, S.~Ren, and J.~Sun,
``Deep residual learning for image recognition,''
in \textit{Proc. CVPR}, 2016, pp.~770--778.

\bibitem{zhang2023biomedclip}
S.~Zhang \textit{et al.},
``BiomedCLIP: a multimodal biomedical foundation model pretrained from
fifteen million scientific image-text pairs,''
\textit{NEJM AI}, 2024;\newline
\textit{arXiv:2303.00915}, 2023.

\bibitem{lin2017focal}
T.-Y. Lin, P.~Goyal, R.~Girshick, K.~He, and P.~Doll\'ar,
``Focal loss for dense object detection,''
in \textit{Proc. ICCV}, 2017, pp.~2980--2988.

\bibitem{lin2024cbam}
Y.-H. Lin, X.-X. Zhang, and L.-Y. Shang,
``Multi-attention enhanced CNN for lumbar disc severity classification
on a balanced RSNA subset,''
\textit{Bioengineering}, 2024.

\bibitem{datascarcity2024}
S.~Woerner and C.~F. Baumgartner,
``Navigating data scarcity using foundation models: a benchmark of
few-shot and zero-shot learning approaches in medical imaging,''
in \textit{Proc. MICCAI 2024 Workshop on Foundation Models for
General Medical AI}, 2024;\newline
\textit{arXiv:2408.08058}.

\bibitem{lu2024radclip}
Z.~Lu \textit{et al.},
``RadCLIP: enhancing radiologic image analysis through contrastive
language-image pre-training,''
\textit{arXiv:2403.09948}, 2024.

\bibitem{yang2026deciphermr}
H.~Yang \textit{et al.},
``Decipher-MR: a 3D MRI vision-language foundation model,''
\textit{npj Digital Medicine}, 2026;\newline
\textit{arXiv:2509.21249}.

\bibitem{blankemeier2026merlin}
L.~Blankemeier \textit{et al.},
``Merlin: a vision-language foundation model for 3D CT,''
\textit{Nature}, 2026.

\bibitem{koleilat2025biomedcoop}
T.~Koleilat \textit{et al.},
``BiomedCoOp: prompt tuning of biomedical vision-language models,''
in \textit{Proc. CVPR}, 2025.

\bibitem{vandergraaf2024spider}
J.~W. van~der~Graaf \textit{et al.},
``Lumbar spine segmentation in MR images: a dataset and a public
benchmark,''
\textit{Scientific Data}, vol.~11, no.~1, 2024.

\bibitem{pfirrmann2001}
C.~W.~A. Pfirrmann, A.~Metzdorf, M.~Zanetti, J.~Hodler, and
N.~Boos,
``Magnetic resonance classification of lumbar intervertebral disc
degeneration,''
\textit{Spine}, vol.~26, no.~17, pp.~1873--1878, 2001.

\bibitem{modic1988}
M.~T. Modic, P.~M. Steinberg, J.~S. Ross, T.~J. Masaryk, and
J.~R. Carter,
``Degenerative disk disease: assessment of changes in vertebral body
marrow with MR imaging,''
\textit{Radiology}, vol.~166, no.~1, pp.~193--199, 1988.

\bibitem{aktan2025}
Y.~Liu, J.~Rashid, S.~M. Boulaaras \textit{et al.},
``A deep learning approach to lumbar spine degenerative classification
on the RSNA 2024 dataset,''
\textit{International Journal of Computational Intelligence Systems},
vol.~18, 2025, doi: 10.1007/s44196-025-01098-7.

\bibitem{hong2025mscan}
J.~Hong \textit{et al.},
``M-SCAN: multi-view scan attention network for lumbar spine
degenerative classification,''
\textit{arXiv:2503.01634}, 2025.

\bibitem{phaphuangwittayakul2026}
A.~Phaphuangwittayakul \textit{et al.},
``Multi-task deep learning for lumbar spine degenerative
classification with external validation,''
\textit{Diseases}, 2026.

\bibitem{nshen2024}
N.~Shen,
``RSNA 2024 lumbar spine degenerative classification baseline,''
GitHub repository, 2024.

\bibitem{kaggle2024}
RSNA, ``RSNA 2024 Lumbar Spine Degenerative Classification
competition leaderboard,''
Kaggle, 2024.

\end{thebibliography}

\end{document}
